\clearpage
\fancyhead[L]{\small EINLEITUNG}
\SourceLeaf{351}
\pdfbookmark[0]{Einleitung zu den Beobachtungen des Magnetischen Vereins von 1836}{einleitung-magnetic-association-1836}
\begin{center}
{\Large EINLEITUNG}\par\medskip
{\small [FÜR DIE ZEITSCHRIFT: RESULTATE AUS DEN BEOBACHTUNGEN DES MAGNETISCHEN VEREINS IM JAHRE 1836. HERAUSGEGEBEN VON C. F. GAUSS UND W. WEBER.]}
\end{center}
\medskip
Die erdmagnetische Kraft äußert sich auf vielfältige Weise. Diese Erscheinungen lassen sich nur durch zahlreiche genaue Beobachtungen ergründen, die an den verschiedensten Punkten der Erdoberfläche fortlaufend angestellt werden. Am meisten bedürfen die \textit{unregelmäßigen Änderungen}, denen diese Kraft, wie wir feststellen, unterliegt, einer streng koordinierten Zusammenarbeit der Beobachter. Es ist hinlänglich bekannt, dass sich die Bestimmungsgrößen der erdmagnetischen Kraft, die Deklination, die Inklination und zweifellos auch die Intensität, fortwährend ändern. Für die Intensität fehlen allerdings noch ausreichende Erfahrungen, da sie erst seit einigen Jahrzehnten untersucht wird. Es gibt säkulare Änderungen, die erst nach längerer Zeit auffallen, im Laufe der Zeit aber sehr beträchtlich werden, und periodische Änderungen, die mit der Jahres- und Tageszeit wechseln. Soweit diese Änderungen regelmäßig verlaufen, ist eine streng koordinierte Zusammenarbeit der Beobachter an verschiedenen Orten zwar höchst wünschenswert, um schneller zu weiteren Erkenntnissen zu gelangen, aber nicht unbedingt nötig. Jeder Beobachter kann auch unabhängig von den anderen nützliche Beiträge leisten.

Anders verhält es sich mit den unregelmäßigen Änderungen, denen man erst in den letzten Jahren größere Aufmerksamkeit zu schenken begonnen hat. Schon vor bei\SourceLeaf{352}nahe hundert Jahren bemerkten \textsc{Hiorter} und \textsc{Celsius}, dass die Magnetnadel bei sichtbarem Nordlicht unregelmäßige und oft sehr große Bewegungen zeigt. Zahlreiche spätere Beobachtungen bestätigten dies. Daraus ließ sich schließen, dass dieselben Kräfte, die ein Nordlicht hervorrufen, zugleich auf die Magnetnadel wirken und dass sich diese Wirkungen über sehr große Entfernungen erstrecken müssen, da Nordlichter gewöhnlich in einem weiten Umkreis sichtbar sind. Eine noch bessere Vorstellung davon, wie weit jene rätselhaften Kräfte wirken, gibt uns die Beobachtung von Herrn \textsc{Arago}. Oft wurden an denselben Tagen, an denen er in Paris starke Störungen der regelmäßigen Bewegung der Magnetnadel beobachtet hatte, an entfernten Orten Nordlichter gesehen, die über dem Horizont von Paris nicht sichtbar waren.

Die Unregelmäßigkeiten in den Erscheinungen des Erdmagnetismus gewannen dadurch ein besonderes Interesse. Gerade auch Herr \textsc{von Humboldt} hatte bei seinen zahlreichen Beobachtungen der täglichen und stündlichen Bewegungen der Magnetnadel bemerkt, wie häufig sie vorkommen. Zwar berechtigten jene Beobachtungen keineswegs dazu, alle unregelmäßigen Bewegungen als gleichzeitig mit Nordlichtern anzusehen. Sie schlossen auch die Möglichkeit nicht aus, dass viele, vielleicht die meisten, allein örtliche Ursachen hatten. Doch ließ sich kaum verkennen, dass dabei nicht selten starke, über weite Entfernungen wirkende Naturkräfte im Spiel sind. Wie diese Kräfte wirken und wie weit sie sich ausbreiten, ist eine lohnende Frage für die Naturforschung, auch wenn ihre Quelle vorerst unbekannt bleibt.

Beiläufig und auf gut Glück gemachte Beobachtungen können uns diesem Ziel nicht näherbringen. Um es zu erreichen, müssen viele solcher Erscheinungen gleichzeitig an vielen Orten im Einzelnen verfolgt und nach Zeit und Größe genau gemessen werden. Dafür sind zuvor genaue Absprachen zwischen Beobachtern nötig, denen geeignete Hilfsmittel zur Verfügung stehen.

Der berühmte Naturforscher, dem unsere Kenntnis des Erdmagnetismus so viel verdankt, hat auch hier zuerst den Weg gebahnt. Gegen Ende des Jahres 1828 errichtete Herr \textsc{von Humboldt} in Berlin eigens für die magnetischen Beobachtungen ein kleines eisenfreies Haus. Er stellte darin einen von \textsc{Gambey} gefertigten Variationskompass auf und vereinbarte mit anderen Be\SourceLeaf{353}sitzern ähnlicher Instrumente an mehreren zum Teil sehr weit entfernten Orten, die magnetische Variation regelmäßig an verabredeten Tagen zu beobachten. Es wurden acht Termine im Jahr, jeder von 44 Stunden. festgesetzt, an denen die magnetische Deklination stündlich aufgezeichnet werden sollte. An einigen Orten beobachtete man in noch kürzeren Abständen, jede halbe Stunde oder alle zwanzig Minuten. Die näheren Bestimmungen findet man im 19. Band von \textsc{Poggendorff}s \textit{Annalen der Physik} auf S.~361. Dort finden sich auch die Beobachtungen, die dieser Vereinbarung gemäß an den Terminen in den Jahren 1829 und 1830 in Berlin, Freiberg, Petersburg, Kasan und Nicolajef angestellt wurden, sowie grafische Darstellungen von dreien davon.

Das hiesige magnetische Observatorium wurde 1833 erbaut. Sein magnetischer Apparat ist ganz anders eingerichtet als die früher benutzten Apparate. Hier wurden die Terminbeobachtungen am 20. und 21. März 1834 zum ersten Mal vollständig angestellt. Entsprechende Beobachtungen wurden nur aus Berlin bekannt. In Göttingen wurde jedoch alle zehn Minuten beobachtet, in Berlin nur stündlich. Gleichwohl zeigten die Berliner Aufzeichnungen mehrere recht beträchtliche Bewegungen, die sich in den Göttinger Beobachtungen wiederfanden. Letztere ließen in den Zwischenzeiten sehr viele andere Bewegungen erkennen, die in Berlin natürlich ganz entgehen mussten. Die Frage, ob ein kleinerer oder größerer Teil der in Göttingen beobachteten Schwankungen nur örtlich gewesen war, blieb daher offen.

Schon der nächste Termin am 4. und 5. Mai brachte jedoch eine Entscheidung. Die Abstände wurden noch weiter verkürzt, und zwar auf fünf Minuten, wodurch die Ergebnisse noch wesentlich deutlicher hervortraten. Entsprechende Beobachtungen mit \textsc{Gambey}schen Apparaten sind weder von diesem noch von einem der späteren Termine bekannt geworden. Herr \textsc{Sartorius} dagegen, der an den Beobachtungen des Märztermins in Göttingen aktiv teilgenommen hatte, hatte sich für eine mehrjährige Reise nach Italien einen Apparat beschafft, der ganz dem Göttinger glich, nur kleiner gebaut war. Damit hatte er den Maitermin in Waltershausen (einem Gut in Bayern, etwa 20 Meilen von Göttingen entfernt) in kurzen Zeitabständen sorgfältig und vollständig beob\SourceLeaf{354}achtet. Es zeigte sich eine wirklich überraschend große Übereinstimmung, sowohl bei den größeren als auch bei fast allen kleineren Schwankungen, die in kurzen Zeitabständen wechselten. Tatsächlich blieb nichts übrig, was man mit Recht örtlichen Ursachen hätte zuschreiben können.

Die drei folgenden Termine im Juni, August und September wurden in Göttingen auf genau dieselbe Weise abgehalten, während die Zahl der auswärtigen Teilnehmer mit ähnlichen oder gleichen Apparaten stetig wuchs. Nachdem Herr Professor \textsc{Encke} die hiesigen Einrichtungen selbst besichtigt hatte, ließ er für Berlin vorläufig einen ähnlichen Apparat in kleinerem Maßstab anfertigen. Herr \textsc{Sartorius} beobachtete mit seinem Apparat an allen Terminen, an denen es die Umstände zuließen (im Juni in Frankfurt, im September in Bramberg im Salzburgischen). In Leipzig, Copenhagen und Braunschweig beobachtete man mit Apparaten, die genau dem hiesigen entsprachen. Das Ergebnis der entsprechenden Beobachtungen ähnelte dem vom Maitermin sehr. Fast alle der zahlreichen in Göttingen beobachteten Schwankungen fanden sich in den Beobachtungen der anderen Orte wieder, wenn auch in anderen Größenverhältnissen, doch in unverkennbarer Übereinstimmung.

Um für dieses bemerkenswerte Ergebnis einen weiteren unwiderlegbaren Beleg zu erhalten, wurden während des damaligen Aufenthalts von Herrn Professor \textsc{Weber} in Leipzig einige besondere gleichzeitige Beobachtungen zwischen Göttingen und Leipzig verabredet. Dafür legte man am 1. und 2. Oktober bestimmte Stunden am Vormittag, mittags und am Abend fest. Diese Beobachtungen wurden von besonders geübten Beobachtern mit größter Sorgfalt durchgeführt. Sie sind in \textsc{Poggendorff}s \textit{Annalen der Physik}, Bd.~33, S.~426, vollständig abgedruckt und durch grafische Darstellungen veranschaulicht.

Damit war sehr deutlich geworden, dass die Erscheinungen in viel kürzeren Zeitabständen verfolgt werden mussten, als Herr \textsc{von Humboldt} sie gewählt hatte. Eine Zeit lang haben wir an den Terminen alle drei Minuten beobachtet, und einige andere Teilnehmer taten dasselbe. Da jedoch ein Teil der auswärtigen Teilnehmer bei Abständen von fünf Minuten blieb, die in den meisten gewöhnlichen Fällen auch ausreichen, haben wir später der Einheitlichkeit halber diese Abstände zur allgemeinen Regel gemacht. Da aber bei so kurzen Zeitabständen die
\SourceLeaf{355}Durchführung der Termine ungleich mühsamer wird als stündliches Aufzeichnen, besonders dort, wo nur wenige Personen sich die Arbeit teilen müssen, so musste. um den Fortbestand des Vereins zu sichern, sowohl die Zahl als auch die Dauer der Termine verringert werden. Seither ist ihre Zahl auf sechs im Jahr und die Dauer jedes Termins auf 24 Stunden festgesetzt. Jedem Haupttermin wurden zwei Nebentermine hinzugefügt. Näheres findet man weiter unten.

Nach dieser Regelung wurden und werden die Beobachtungen in Göttingen und an immer mehr anderen Orten ohne Unterbrechung fortgesetzt. Apparate, die dem Göttinger gleichen oder ähneln, gibt es in Altona, Augsburg, Berlin, Bonn, Braunschweig, Breda, Breslau, Cassel, Copenhagen, Dublin, Freiberg, Göttingen, Greenwich, Halle, Kasan, Krakau, Leipzig, Mailand, Marburg, München, Neapel, Petersburg und Upsala. Von acht dieser Orte sind uns bisher keine Beobachtungen bekannt geworden. An einigen der übrigen Orte ist die Teilnahme an den Beobachtungen aufgrund äußerer Umstände noch nicht regelmäßig und ohne Unterbrechung.

Aus der ersten Zeit dieses Zusammenschlusses wurden einige Termine in \textsc{Schumacher}s \textit{Astronomischen Nachrichten} und \textsc{Poggendorff}s \textit{Annalen der Physik} in grafischer Form veröffentlicht. Da die Teilnahme inzwischen so stark gewachsen war, schien es an der Zeit, für eine regelmäßige Veröffentlichung zu sorgen, um diese Fülle ergiebiger Tatsachen allen zugänglich zu machen, die sich für die Naturforschung interessieren. Was wir jetzt vorlegen, kann als der erste Jahrgang. seitdem der Verein einen gewissen Umfang erreicht hat, gelten. Vom Jahr 1837 an werden die Ergebnisse jedes Termins regelmäßig veröffentlicht, sobald sie in ausreichender Vollständigkeit vorliegen.

Die Beobachtungen und ihre grafischen Darstellungen sollen von Erläuterungen und Bemerkungen begleitet werden, die unmittelbar dazu gehören, und außerdem von weiteren Aufsätzen. Darin sollen Themen aus dem weiten Gebiet des Erdmagnetismus, die dafür verwendeten Instrumente, ihre Justierung und Handhabung und. die vielfältigen Anwendungen dieser Instrumente Platz finden.

\SourceLeaf{356}Zum unmittelbaren Gegenstand der Arbeit unseres Vereins, den Änderungen der magnetischen Deklination, sei noch eine Bemerkung erlaubt. Wenn, wie außer Zweifel steht, die beiden anderen Bestimmungsgrößen der erdmagnetischen Kraft, die Inklination und die Intensität, ähnlichen Änderungen unterliegen, kann man fragen, warum dem ersten Element vorzugsweise oder vorerst ausschließlich so große Sorgfalt gewidmet wird.

Die Kenntnis der Änderungen und Störungen der magnetischen Deklination ist in der Tat von sehr großem praktischem Interesse. Für den Seefahrer, den Geodäten und den Markscheider muss es von größter Bedeutung sein, zu wissen, wie häufig und wie stark eines ihrer wichtigsten Arbeitsmittel unvermeidlich gestört wird, und sei es nur, um abzuschätzen, wie weit sie ihm vertrauen dürfen. Für die beiden zuletzt genannten Anwendungen der Bussole, in der praktischen Geometrie über und unter Tage, können diese Untersuchungen in Zukunft sogar noch viel nützlicher werden. Wenn einmal feststeht, dass die mit der Zeit wechselnden unregelmäßigen Störungen nie oder nur äußerst selten rein örtlich sind, sondern immer oder fast immer über weite Gebiete ganz gleichzeitig und in fast gleicher Stärke auftreten, lassen sie sich fast vollständig unschädlich machen. Der Geodät und der Markscheider brauchen nur bei allen Arbeiten mit der Bussole genau auf die Uhrzeit zu achten und gleichzeitig an einem anderen, nicht allzu weit entfernten Ort Beobachtungen anstellen zu lassen. Durch Vergleich lassen sich diese Störungen ebenso ausschalten, wie reisende Beobachter ihre barometrischen Höhenbestimmungen durch den Vergleich mit Barometerbeobachtungen an einem festen Ort von den unregelmäßigen Schwankungen des Barometerstands unabhängig machen. Es versteht sich von selbst, dass hier keine Störungen gemeint sind, denen die Bussole in eisenhaltigen Gruben ausgesetzt ist.

Der Grund dafür, dass die Deklination den anderen Bestimmungsgrößen des Erdmagnetismus vorgezogen wird, liegt dennoch weniger in diesen praktischen Erwägungen als im gegenwärtigen Stand der Hilfsmittel. Für den Naturforscher hat das Aufsuchen von Gesetzen in den Naturerscheinungen seinen Zweck und Wert in sich selbst, und es liegt ein besonderer Zauber darin, Maß und Harmonie im anschei\SourceLeaf{357}nend völlig Regellosen zu erkennen. Beim Verfolgen des wunderbaren Spiels der ständig wechselnden Änderungen der Deklination lassen die heute verwendeten Apparate an Sicherheit, Genauigkeit und Leichtigkeit der Beobachtungen nichts zu wünschen übrig. Von den bisherigen Beobachtungsmitteln für die beiden anderen Bestimmungsgrößen lässt sich das nicht sagen. Zurzeit ist es deshalb noch zu früh, sie in groß angelegte Untersuchungen einzubeziehen. Sobald aber die Beobachtungsmittel so weit verbessert sind, dass wir die Änderungen, insbesondere die rasch wechselnden Änderungen, bei den anderen Bestimmungsgrößen des Erdmagnetismus sicher erkennen, leicht verfolgen und genau messen können, werden diese Änderungen denselben Anspruch auf die vereinten Anstrengungen der Naturforscher haben wie die Änderungen der Deklination. Man darf hoffen, dass dieser Zeitpunkt nicht mehr allzu fern ist.

